\section{Changes to the Task 1 Metamodels}
\label{app:metamodels}

Writing the transformation changed the Task~1 metamodels in three ways. A
third metamodel, the pinned inputs, was added
(Section~\ref{app:pinned}). The target metamodel gained the values the
transformation derives that Task~1 had not yet recorded, and its
generated-file classes now cover applications with several processes
(Section~\ref{app:target}). Four grammars were added for the pinned inputs
(Section~\ref{app:grammars}). The source metamodel did not change.

\subsection{The Pinned Inputs Metamodel}
\label{app:pinned}

The package \texttt{pinnedinputs} contains the documents the resolution reads
beside the source models. They are a separate package because they are not
part of what an author writes. The source metamodel's structure is compared
with the production implementation's schemas through a committed descriptor,
and adding these classes to it would have added classes those schemas do not
contain. No class of the target metamodel refers to them: the resolved model
records the values the transformation reads from them, and their digests.
Table~\ref{tab:pinned} lists the classes.

\begin{table}[!htbp]
\centering
\small
\begin{tabular}{>{\raggedright\arraybackslash}p{0.2\linewidth}>{\raggedright\arraybackslash}p{0.36\linewidth}>{\raggedright\arraybackslash}p{0.36\linewidth}}
\hline
\textbf{Class} & \textbf{Features} & \textbf{Read by} \\
\hline
\texttt{NodeContract} & cluster, label prefix, \texttt{nodes} & eligibility \\
\texttt{Node} & name (ID), site, architecture, \texttt{allocatable}, capabilities, \texttt{gpus}, \texttt{disks} & eligibility \\
\texttt{Allocatable} & cpu, memory, as quantities such as \texttt{53500m} & eligibility \\
\texttt{Gpu} & vendor, model, class, memory in MiB & GPU requests \\
\texttt{Disk} & medium, usable GiB & disk requests \\
\texttt{ImagesLock} & name, \texttt{images} & every image \\
\texttt{LockedImage} & alias (ID), repository, digest, user, group & images, sidecars, backups, migrations \\
\texttt{ClusterState} & cluster, capture time, \texttt{bindings}, \texttt{placements} & bound machines \\
\texttt{Binding} & claim, node & bound machines \\
\texttt{CurrentPlacement} & process, node & not yet read \\
\texttt{MigrationProof} & \texttt{applications} & migrations \\
\texttt{ProvenApplication} & application id, tested schema version (optional), whether the release holds a change that cannot run in a transaction & migrations \\
\texttt{AssetFile} & project, path, text & configuration files \\
\hline
\end{tabular}
\caption{The Classes of the Pinned Inputs Metamodel}
\label{tab:pinned}
\end{table}

The disk medium of a node is a string, not the source metamodel's
\texttt{Media} enumeration. The node contract may list a medium that no process
may request: three machines boot from SD cards, and \texttt{sdcard} is not a
\texttt{Media} literal. \texttt{AssetFile} is not parsed from YAML. The pipeline
creates one per configuration file that a project names, from the file beside
the project document.

\subsection{Changes to the Target Metamodel}
\label{app:target}

Table~\ref{tab:targetvalues} lists the values added so that the transformation
can record what it derives, and Table~\ref{tab:targetfiles} the changes to the
generated-file classes, which let one file describe an application with several
processes and cover every kind of file the specification's generated trees
hold. The Task~1 model of \texttt{minimal} needed neither.

\begin{table}[!htbp]
\centering
\small
\begin{tabular}{>{\raggedright\arraybackslash}p{0.24\linewidth}>{\raggedright\arraybackslash}p{0.33\linewidth}>{\raggedright\arraybackslash}p{0.35\linewidth}}
\hline
\textbf{Class} & \textbf{Change} & \textbf{Reason} \\
\hline
\texttt{StartupProbe} & added \texttt{timeout} & The specification gives the startup probe the platform's probe timeout. \\
\texttt{ResolvedGrant} & \texttt{path} and \texttt{access} optional; added \texttt{engine} and \texttt{paths} & A transit or database secret has no single path. It is recorded as the Vault paths its policy covers. \\
\texttt{PolicyPath} & new: \texttt{path}, \texttt{allows} & One path of a secret's policy and the operations allowed on it. \\
\texttt{ResolvedVolume} & added \texttt{mountAt} & The workload template needs the path a volume is mounted at. \\
\texttt{ResolvedProcess} & added \texttt{lifecycle} & Whether the process serves, runs once or prepares; the revision's element names it. \\
\texttt{ResolvedScrape} & new: \texttt{process}, \texttt{surface}, \texttt{path}, \texttt{interval}, \texttt{timeout} & How the metrics stack scrapes the application, which the revision's element names. \\
\texttt{ReleaseGate} & added \texttt{endpoint} & The address every Canary's webhooks call. \\
\texttt{ResolvedExposure} & added \texttt{listener}, \texttt{certificates}, \texttt{proxyApplication}, \texttt{proxyNamespace} & The entry point and certificate source of the tier, and the proxy whose middleware the routes name. \\
\texttt{PolicyPeer} & added \texttt{rule} & Which rule admitted the peer, so two rules admitting one peer stay two entries, as in the production implementation. \\
\hline
\end{tabular}
\caption{Values Added to the Target Metamodel}
\label{tab:targetvalues}
\end{table}

\begin{table}[!htbp]
\centering
\small
\begin{tabular}{>{\raggedright\arraybackslash}p{0.24\linewidth}>{\raggedright\arraybackslash}p{0.33\linewidth}>{\raggedright\arraybackslash}p{0.35\linewidth}}
\hline
\textbf{Class} & \textbf{Change} & \textbf{Reason} \\
\hline
\texttt{WorkloadFile}, \texttt{SwitchoverFile}, \texttt{IdentityFile} & \texttt{process} replaced by \texttt{processes} (one or more) & One file per application holds the resources of all its processes, as in the generated tree of the specification. \\
\texttt{NetworkPolicyFile} & \texttt{process}, \texttt{ingress} and \texttt{egress} replaced by \texttt{policies} & One policy file per application holds one policy per process. The project's default-deny file holds none. \\
\texttt{ProcessPolicy} & new: \texttt{process}, \texttt{ingress}, \texttt{egress} & The peers that may reach one process, and those it may reach. \\
\texttt{MonitorFile} & added \texttt{timeout} and \texttt{kind} & A blue/green process is scraped through its pods, any other through its service. \\
\texttt{ServiceFile}, \texttt{ClaimFile}, \texttt{ConfigMapFile}, \texttt{BackupFile} & new, each with \texttt{processes} & The services, volume claims, configuration files and backups of an application's processes. \\
\texttt{SecretFile} & \texttt{process} and \texttt{grants} replaced by \texttt{application} and \texttt{holders} & One file per application holds the secrets of every identity in it, including a backup identity. \\
\texttt{SecretHolder} & new: \texttt{identity}, \texttt{component}, \texttt{grants} & One identity and the secrets synced for it. \\
\texttt{ConnectionFile} & new: \texttt{namespace}, \texttt{address} & The secret operator's connection to the secret store. \\
\texttt{VaultPolicyFile}, \texttt{VaultRoleFile} & new, with \texttt{namespace} and \texttt{identity} & The Vault policy and role of one identity, written as JSON. \\
\hline
\end{tabular}
\caption{Changes to the Generated-File Classes}
\label{tab:targetfiles}
\end{table}

The hand-written model of \texttt{minimal} follows these changes. Its revision
is now the production implementation's, and its generated files refer to the
process through the new features.

\subsection{Grammars for the Pinned Inputs}
\label{app:grammars}

Four grammars were added beside the project and platform grammars:
\texttt{NodeContract}, \texttt{ImagesLock}, \texttt{ClusterState} and
\texttt{MigrationProof}. Each inherits the project grammar, so it reuses its
terminals, its indentation tokens and its scalar rule, and adds only the rules
for its own document and the keywords it introduces. Listing~\ref{lst:grammar}
shows the smallest.

\begin{lstlisting}[style=report, caption={The Grammar of the Cluster-State Snapshot}, label={lst:grammar}]
grammar dev.jorisjonkers.deploykit.emf.syntax.ClusterState
	with dev.jorisjonkers.deploykit.emf.syntax.ProjectIntent
import "https://jorisjonkers.dev/deploy-kit/pinned-inputs/1"

ClusterState returns ClusterState:
	(('apiVersion' ':' apiVersion=Text)
	& ('kind' ':' kind=Text)
	& ('schemaVersion' ':' schemaVersion=Text)
	& ('cluster' ':' cluster=Text)
	& ('capturedAt' ':' capturedAt=Text)
	& ('bindings' ':' BEGIN (DASH BEGIN bindings+=Binding END)* END)
	& ('placements' ':' BEGIN (DASH BEGIN placements+=CurrentPlacement END)* END));
\end{lstlisting}

As with the platform document, the file name selects the grammar:
\texttt{node-contract.yml}, \texttt{images.lock.yml},
\texttt{cluster-state.yml} and \texttt{migration-proof.yml}. One input file
changed. The scalar terminal of the project grammar cannot contain a colon, so
the snapshot's capture time is now quoted, as digests and URLs already were.
The production implementation reads the same value either way.
